\section[Simulation Methodology]{Multiphysics Simulation Methodology}
\label{sec:pd_simulation}

The device performance metrics established in \cref{sec:pd_metrics} are evaluated
through a coupled optical-electrical simulation chain whose output is anchored to
the experimental benchmark of Shi~\textit{et al.}~\cite{Shi2024}. This section
specifies each solver, its numerical settings, the data flow connecting the
stages, and the reduced-order checks used to compare the solver outputs with the
geometry-derived bandwidth and noise models implemented in MATLAB.

\subsection[Optical Simulation: FDTD]{Optical Simulation: Three-Dimensional FDTD}
\label{subsec:pd_fdtd}

The optical simulation solves the time-domain Maxwell's equations of
\cref{eq:maxwell_curl_e,eq:maxwell_curl_h} throughout a domain encompassing
the silicon input waveguide, the \SI{40}{\micro\metre} adiabatic taper, the
complete germanium absorber, and the surrounding SiO$_2$ cladding. The solver
(Lumerical FDTD Solutions 2020~R2) discretises the electromagnetic field on a
Yee staggered grid using the second-order Courant-stable leapfrog scheme.
Computational photonics methodology for FDTD-based device simulation, including
the algorithmic implementation of the Yee algorithm and absorbing boundary
conditions, is detailed in the pedagogical treatment of Wartak~\cite{Wartak2013}.

The complex permittivity is sampled at the material-dependent optical frequency
through tabulated complex refractive indices. The auxiliary MATLAB response
model uses $\alpha\approx\SI{7e3}{\per\centi\metre}$ at
\SI{1310}{\nano\metre} for analytical responsivity estimates, while the FDTD
script evaluates absorption directly from the Palik material model.

\paragraph{FDTD source field and the $\varepsilon''$-based generation rate.}
Following the convention of Hamadou~\textit{et al.}~\cite{Hamadou2023}, the
monochromatic optical source in the FDTD domain is a Gaussian-modulated pulse:
\begin{equation}
    \tilde{\mathbf{E}}(\mathbf{r},t)
    = E_0\exp\!\left[-\frac{(t-t_0)^2}{2\Delta t^2}\right]
      \sin\!\bigl[-\omega(t-t_0)\bigr]\,\hat{\mathbf{x}},
    \label{eq:fdtd_source}
\end{equation}
where $t_0$ is the temporal offset and $\Delta t$ is related to the
pulse FWHM by $\Delta t=\mathrm{FWHM}/(2\sqrt{\ln 2})$.
After the time-domain solution is Fourier-transformed, the volumetric
optical generation rate extracted from the monochromatic field is:
\begin{equation}
    G(\mathbf{r},\omega)
    = \frac{\pi\,\varepsilon_0\,\varepsilon''_r(\omega)}{h}
      \,|\tilde{\mathbf{E}}(\mathbf{r},\omega)|^2,
    \label{eq:gen_rate_eps}
\end{equation}
where $\varepsilon''_r(\omega)=2n\kappa$ is the imaginary part of the relative
permittivity and $h$ is Planck's constant~\cite{Hamadou2023}.
\Cref{eq:gen_rate_eps} is equivalent to \cref{eq:gopt_poynting_local} under
the plane-wave approximation, and reduces to \cref{eq:gen_rate} in the
uniform-field limit; the FDTD solver evaluates it pointwise over the
three-dimensional Ge volume before transferring the resulting
$G(\mathbf{r})$ array to the CHARGE drift-diffusion solver.
Frequency-domain power monitors record the input-normalised transmission before
and after the Ge absorber, from which the script computes:
\begin{equation}
    \mathcal{A}(\omega)
    = \max\!\left[\frac{T_{\mathrm{ref}}(\omega)-T_{\mathrm{after}}(\omega)}
                        {T_{\mathrm{ref}}(\omega)},\,0\right].
    \label{eq:fdtd_abs}
\end{equation}
A \SI{20}{\nano\metre} isotropic mesh override within the germanium region
resolves the exponential generation-rate gradient.

\paragraph{FDTD numerical stability and the 3D Courant condition.}
The PML implementation employs eight absorbing layers on each domain boundary,
matching the setting in \texttt{ge\_pd\_fdtd\_oband\_ushaped.lsf}. The temporal stability
of the Yee algorithm requires the 3D Courant-Friedrichs-Lewy condition:
\begin{equation}
    \Delta t < \frac{\Delta}{c\sqrt{3}},
    \label{eq:cfl_3d}
\end{equation}
so the time step is set to $\Delta t=0.99\,\Delta/(c\sqrt{3})$. The spatial
mesh size is independently constrained by the sampling requirement
$\Delta<\lambda_{\min}/(10\sqrt{\varepsilon_{\max}})$, giving
$\Delta<\SI{32}{\nano\metre}$ with a safety margin of 1.6 at the
\SI{20}{\nano\metre} override.

\paragraph{FDTD electromagnetic solution and the square-law photocurrent.}
A physically essential clarification connecting the FDTD solution to the
photocurrent calculation concerns what the power monitors extract. The
frequency-domain power recorded is the time-averaged Poynting vector
$\mathbf{S}(\mathbf{r},\omega)=\tfrac{1}{2}\mathrm{Re}[\tilde{\mathbf{E}}\times\tilde{\mathbf{H}}^*]$,
which is quadratic in the field amplitudes, consistent with the square-law
mechanism of \cref{eq:square_law}. The optical generation rate
$G_{\mathrm{opt}}(\mathbf{r})=\mathcal{P}_{\mathrm{abs}}(\mathbf{r})/\hbar\omega$
of \cref{eq:gen_rate} is thus already expressed in the power-quadratic form that
maps directly to the photocurrent via $I_{\mathrm{ph}}=q\int_\Omega G_{\mathrm{opt}}\,\mathrm{d}V$~\cite{Pilori2019}.

\paragraph{Waveguide insertion loss and chip-level responsivity correction.}
The saturation of $\mathcal{A}(L)$ identified in the length sweep is not
smooth when waveguide routing loss is accounted for. For a Si ridge waveguide
with propagation loss $\gamma_{\mathrm{wg}}\approx\SI{3.6}{\decibel\per\centi\metre}$
at \SI{1310}{\nano\metre}, a routing length of $\ell_{\mathrm{wg}}=\SI{1}{\milli\metre}$
introduces an approximately \SI{0.36}{\decibel} correction equivalent to a
\SI{4}{\percent} responsivity reduction relative to the taper-input
value~\cite{Wang2024}. The FDTD domain of the present work terminates at the
taper entry; any comparison with chip-level measurements must apply the factor
$10^{-\gamma_{\mathrm{wg}}\ell_{\mathrm{wg}}/10}$ to \cref{eq:responsivity}.

\begin{figure}[H]
    \centering
    \missingfigure[figwidth=0.80\linewidth]{
        Simulation flowchart connecting FDTD, CHARGE, and MATLAB post-processing.
        Left: FDTD inputs and outputs.
        Centre: CHARGE inputs and outputs.
        Right: CML dataset, compact model, DSP model, and thesis figures.
        Arrows indicate data flow; the key coupling step is the
        $G_{\mathrm{opt}}$ file transfer.}
    \caption{Multi-physics simulation workflow. The FDTD-to-CHARGE coupling
    proceeds via the volumetric generation rate $G_{\mathrm{opt}}(\mathbf{r})$
    of \cref{eq:gen_rate}. Device-level and system-level MATLAB scripts then
    convert the solver outputs into figures, compact-model data, and
    geometry-derived bandwidth estimates.}
    \label{fig:sim_flowchart}
\end{figure}

\subsection[Electrical Simulation: CHARGE]{Electrical Simulation: Two-Dimensional Drift-Diffusion}
\label{subsec:pd_charge}

Carrier transport is simulated in Lumerical DEVICE CHARGE under
two-dimensional X-normal steady-state conditions with the norm length set to
the Ge absorber length $L=\SI{8}{\micro\metre}$. The solver evaluates the
vertical cross-section and scales currents by this longitudinal length, solving
the PDE system of
\cref{eq:poisson,eq:Jn,eq:Jp,eq:cont_n,eq:cont_p} on an adaptive
finite-element mesh. Two complementary carrier mobility models are employed.
The \emph{Caughey-Thomas} model accounts for impurity scattering through a
doping-concentration-dependent low-field mobility~\cite{Hamadou2023}:
\begin{equation}
    \mu_{n,p}(N_A+N_D)
    = \mu^{\min}_{n,p}
      + \frac{\mu^L_{n,p} - \mu^{\min}_{n,p}}
             {1 + \left(\dfrac{N_A+N_D}{N^0_{n,p}}\right)^{\!\alpha_{n,p}}},
    \label{eq:sim_caughey_thomas}
\end{equation}
where $\mu^L_{n,p}$ is the lattice-scattering-limited mobility, $\mu^{\min}_{n,p}$
the ionised-impurity saturation floor, and $N^0_{n,p}$, $\alpha_{n,p}$ fitting
parameters. High-field velocity saturation is then superimposed through a
monotonic saturation model~\cite{Hamadou2023}:
\begin{equation}
    \mu'_{n,p}(|\mathbf{E}|)
    = \frac{\mu_{n,p}}
           {\left[1
             + \left(\dfrac{\mu_{n,p}|\mathbf{E}|}{v^S_{n,p}}
               \right)^{\!\beta_{n,p}}
           \right]^{\!1/\beta_{n,p}}},
    \label{eq:highfield_mob}
\end{equation}
yielding the field-dependent diffusivity $D_{n,p}=\mu'_{n,p}k_BT/q$ from
\cref{eq:einstein}. \Cref{eq:highfield_mob} is equivalent in functional form
to the Canali velocity-saturation model~\cite{Alasio2024}:
\begin{equation}
    \mu_{n,p}(|\mathbf{E}|)
    = \frac{(\alpha_{\mathrm{C}}+1)\,\mu_{n,p}^{(0)}}
      {\left[\alpha_{\mathrm{C}}
        + \!\left(1 + \!\left(\frac{(\alpha_{\mathrm{C}}+1)\,
          \mu_{n,p}^{(0)}\,|\mathbf{E}|}{v_{\mathrm{d,sat}}}\right)^{\!\beta}
        \right)^{1/\beta}\right]},
    \label{eq:canali_full}
\end{equation}
which reduces to the saturation model of \cref{eq:mobility_sat} with $\beta=2$
and $\alpha_{\mathrm{C}}=0$. The low-field mobilities adopted are
$\mu_n^{(0)}=\SI{3900}{\centi\metre\squared\per\volt\per\second}$ and
$\mu_p^{(0)}=\SI{1900}{\centi\metre\squared\per\volt\per\second}$; the
saturation velocities
$v_{\mathrm{d,sat},e}=\SI{6e6}{\centi\metre\per\second}$ and
$v_{\mathrm{d,sat},h}=\SI{4e6}{\centi\metre\per\second}$ enter
\cref{eq:canali_full} as asymptotic limits. A $\pm20\%$ variation in
$\mu_p^{(0)}$ changes $f_{\mathrm{tt}}$ by less than \SI{1}{\percent},
confirming the robustness of the transit-time estimate in \cref{eq:bw_tt}
to low-field mobility uncertainty~\cite{Alasio2024}.

\paragraph{Junction capacitance and the microhole filling factor.}
For a planar junction the parallel-plate capacitance follows \cref{eq:cj}.
If photon-trapping microholes of total area $A_h$ are etched through the
junction region, the effective junction area is reduced by the filling
factor~\cite{Hamadou2023}:
\begin{equation}
    \rho_f = \frac{A_h}{A_j},
    \quad
    C_{j,\mathrm{mh}}
    = \frac{\varepsilon_0\varepsilon_r\,A_j}{W_i}\,(1-\rho_f),
    \label{eq:cj_filling}
\end{equation}
so that microholes simultaneously improve absorption efficiency and reduce
capacitance, at the cost of increased surface recombination along the
etched sidewalls through \cref{eq:jdark_surf}. The present Ge-on-Si design
does not employ microholes; \cref{eq:cj_filling} is retained as a reference
for future microstructure extensions.

\begin{equation}
    \mu_{n,p}(|\mathbf{E}|) = \frac{\mu_{n,p}^{(0)}}
    {\left[1 + \left(\mu_{n,p}^{(0)}|\mathbf{E}|/v_{\mathrm{d,sat},{n,h}}\right)^{\beta}\right]^{1/\beta}},
    \label{eq:mobility_sat}
\end{equation}
SRH recombination is activated in the germanium absorber with calibrated
lifetimes $\tau_n=\tau_p=\SI{1}{\nano\second}$ in the current repository
configuration. Ohmic boundary conditions are imposed at the cathode
($V=0$) and the anode, which is swept from \SI{-2.5}{\volt} to
\SI{0.5}{\volt} over 51 points. The dark sweep is run first; the FDTD
generation-rate dataset is then imported and the illuminated sweep is repeated.
The same CHARGE setup also runs an SSAC frequency sweep from \SI{1}{\mega\hertz}
to \SI{200}{\giga\hertz}. All performance metrics are
extracted directly from the CHARGE solution: $\mathcal{R}$ from
\cref{eq:responsivity}; $\eta_{\mathrm{int}}$ and $\eta_{\mathrm{ext}}$ from
\cref{eq:qe_full}; $C_j$ from the small-signal AC analysis; and
$f_{3\,\mathrm{dB}}$ from \cref{eq:bw_total} using the extracted $C_j$ and
modelled $R_s$ from \cref{eq:rs_anode}.

\subsection[System-Level Post-Processing: MATLAB]{System-Level Post-Processing: MATLAB}
\label{subsec:pd_matlab}

The MATLAB post-processing pipeline
(\texttt{system-level/ge\_pd\_dsp\_model.m} and
\texttt{device-level/ge\_pd\_oband\_ushaped\_postprocess.m}) consumes the CHARGE
output to compute the system-level performance projections. The key computations
are: (i) the combined transfer function from \cref{eq:bw_total,eq:H_pd_pk}
applied to the extracted $C_j$, modelled $R_s$, and the benchmark-calibrated
peaking parameters $(L_p,C_p)$; (ii) the sensitivity estimate from
\cref{eq:sens_full} at the projected bandwidth; and (iii) the parametric design
sweeps (absorber length and $W_i$ thickness) that produce the trade-off tables
and design-space visualisations.

\subsection[Validation Approach]{Validation Approach and Reference Design}
\label{subsec:pd_validation}

Validation is performed by comparing three levels of the same implemented
pipeline. First, the FDTD absorption monitor is checked against the expected
Beer-Lambert trend over the Ge-length sweep generated by
\texttt{ge\_pd\_fdtd\_oband\_ushaped\_sweeps.lsf}. Second, the CHARGE dark and
illuminated sweeps are reduced to $I_d$, $I_{\mathrm{ph}}$, $\mathcal{R}$,
EQE, and IQE by \texttt{ge\_pd\_oband\_ushaped\_postprocess.m}. Third, the
geometry-derived RC and transit-time model is evaluated from the same dimensions
and compared with the optional SSAC response exported by CHARGE. This validation
does not claim an implemented physics-informed neural-network surrogate; such a
surrogate remains a future extension rather than part of the present repository.


The carrier transport physics that governs the electrical response of
the Ge absorber---including drift-diffusion transport, velocity
saturation, Caughey-Thomas impurity scattering, SRH recombination,
and the bandwidth-absorption trade-off---has been established in
\cref{sec:pd_fundamentals,sec:pd_metrics,sec:pd_material}. The CHARGE
solver evaluates the complete PDE system of
\cref{eq:poisson,eq:Jn,eq:Jp,eq:cont_n,eq:cont_p} self-consistently,
using the Caughey-Thomas doping-dependent mobility of
\cref{eq:sim_caughey_thomas} and the high-field saturation model of
\cref{eq:highfield_mob}. The Ge SRH lifetime is configured as
$\tau_n=\tau_p=\SI{1}{\nano\second}$, consistent with the
dislocation-limited material discussed in
\cref{subsec:pd_ge_tradeoffs}. The transit-time bandwidth of
\cref{eq:bw_tt}, RC bandwidth of \cref{eq:bw_rc}, and their quadrature
combination in \cref{eq:bw_combined} are evaluated from the
geometry-derived capacitance and U-shaped anode resistance model of
\cref{sec:pd_design}; no re-derivation of these expressions is
necessary here.